% Authority: exact-scope leaf 24; full PDF page 103; printed p.86. Translation of direct transcription.
% U:p086.u01 | continuation of p085.u04
\noindent
simplified by means of a few considerations that may not immediately
present themselves to the mind. We begin with the demonstration of the first
of the theorems cited.
% U:p086.u02 | paragraph-start
\par\indent
Let us denote by $p$ a prime number $4n+1$, by $b$ a prime number
$4n+3$, such that $\leg{b}{p}=1$, and consider the equation:
\begin{equation*}\tag{$\alpha'$}
t^2-bu^2=ps^2.
\end{equation*}
% U:p086.u03 | continuation
By the theorem already cited (\emph{Theory of Numbers}, no. 27), the necessary
conditions for this equation to be possible are these: $\leg{b}{p}=1$ and
$\leg{p}{b}=1$. Now, of these conditions the first holds by hypothesis, and as
for the second, it follows from the theorem known as the law of reciprocity
that it always holds when the first is satisfied. The equation is therefore
soluble. It is evident that one may suppose the numbers $t$, $u$, $s$ that
satisfy it to be relatively prime when compared two at a time. The equation being
in this state, one of these three numbers will be even, the other two odd, and
one readily ascertains that $s$ cannot be even. The number $s$ being odd,
of the numbers $t$, $u$, one will be even, the other odd, and there would be two cases to
examine according as $t$ is even or odd. But although these two cases can
be treated by the same analysis, it is simpler to consider only
one of them, that of $t$ odd for example, and to show that the other
can easily be reduced to it. Let us therefore show that it is always permissible to
suppose $t$ odd in the equation $t^2-bu^2=ps^2$. If $t$ were even in
the equation $(\alpha')$, one would deduce from the numbers $t$, $u$, by the means indicated in the note
relating to the beginning of paragraph 3 of the preceding memoir, other
numbers $t'$ (odd), $u'$ (even) such that $t'^2-bu'^2=ps^2$. We may therefore
regard $t$ as odd in all that follows.
% U:p086.u04 | paragraph-start
\par\indent
Let us decompose the numbers $t$ and $u$ into their simple factors, putting:
\[
t=ll'l''\ldots,\quad u=2^\nu kk'k''\ldots,
\]
$l$, $l'$, $l''$, $\ldots$ and $k$, $k'$, $k''$, $\ldots$ denoting odd prime numbers. Mere
inspection of the equation $(\alpha')$ gives:
\[
\leg{p}{k}=1,\quad\leg{p}{k'}=1,\quad\leg{p}{k''}=1,\quad\ldots.
\]
One will then have, on applying the law of reciprocity, $p$ being a prime
number $4n+1$:
\[
\leg{k}{p}=1,\quad\leg{k'}{p}=1,\quad\leg{k''}{p}=1,\quad\ldots.
\]
